% !TEX root = main.tex
\documentclass[main.tex]{subfiles}
\begin{document}

\appendix

\section{Derivatives of $\energy$ and $(\dual,\combweight)$}
\label{sec:appendix-prelim}

Both our objectives $\centroiddual$ and $\jointdual$ are the plain power diagram energy $\energy$ in \eqref{eq:power-diagram} evaluated at the shifted weighted sites $(\dual, \combweight)$ of \eqref{eq:shifted-site}.
We summarize here the ingredients shared by the next two sections: the Jacobian of \eqref{eq:shifted-site} and the derivatives of \eqref{eq:power-diagram}.

The first-order Jacobian of the reparametrization \eqref{eq:shifted-site} are
\begin{equation}
  \partial_{\shift_i} \dual_i = \Id , \qquad
  \partial_{\shift_i} \combweight_i = 2\,(\dual_i - \tcentroid_i) , \qquad
  \partial_{\weight_i} \combweight_i = 1 ,
  \label{eq:appendix-jac}
\end{equation}
and its only nonzero second-order part is $\smash{\partial^2_{\shift_i} \combweight_i = 2\,\Id}$.

As detailed in \cite[\S2.4--2.5]{deGoes2014Thesis}, the first derivatives of the energy $\energy$ in \eqref{eq:power-diagram} are the standard weight and site gradients:
\begin{equation}
  \partial_{\dual_i} \energy = 2 \int_{\cell_i} (\dual_i - \xx) \, d\xx ,
  \qquad
  \partial_{\combweight_i} \energy = -\cellvol_i ,
  \label{eq:appendix-energy-deriv}
\end{equation}
and its second derivatives are facet integrals with off-diagonals
\begin{subequations}
  \label{eq:appendix-energy-hess-off}
  \begin{empheq}[left={\empheqlbrace\,}]{align}
    \partial^2_{\dual_i \dual_j} \energy
      &= \frac{2}{\norm{\dual_i - \dual_j}} \int_{\face_{ij}}\! (\xx - \dual_i)(\xx - \dual_j)\tr \, dA , \label{eq:appendix-energy-hess-pp} \\
    \partial^2_{\dual_i \combweight_j} \energy
      &= \frac{1}{\norm{\dual_i - \dual_j}} \int_{\face_{ij}}\! (\xx - \dual_i) \, dA , \label{eq:appendix-energy-hess-pl} \\
    \partial^2_{\combweight_i \combweight_j} \energy
      &= \frac{1}{2}\frac{\area_{ij}}{\norm{\dual_i - \dual_j}} , \label{eq:appendix-energy-hess-ll}
  \end{empheq}
\end{subequations}
and diagonals
\begin{subequations}
  \label{eq:appendix-energy-hess-diag}
  \begin{empheq}[left={\empheqlbrace\,}]{align}
    \partial^2_{\dual_i \dual_i} \energy
      &= 2\,\cellvol_i\,\Id - \!\sum_{j \sim i} \frac{2}{\norm{\dual_i - \dual_j}} \int_{\face_{ij}}\! (\xx - \dual_i)(\xx - \dual_i)\tr \, dA , \label{eq:appendix-energy-hess-pp-diag} \\
    \partial^2_{\dual_i \combweight_i} \energy
      &= -\!\sum_{j \sim i} \frac{1}{\norm{\dual_i - \dual_j}} \int_{\face_{ij}}\! (\xx - \dual_i) \, dA , \label{eq:appendix-energy-hess-pl-diag} \\
    \partial^2_{\combweight_i \combweight_i} \energy
      &= -\!\sum_{j \sim i} \frac{1}{2}\frac{\area_{ij}}{\norm{\dual_i - \dual_j}} , \label{eq:appendix-energy-hess-ll-diag}
  \end{empheq}
\end{subequations}
where the site--site diagonal carries an extra interior term $2\,\cellvol_i\,\Id$ from the explicit dependence of $\partial_{\dual_i}\energy$ on $\dual_i$, and otherwise each diagonal is the negated neighbor sum of its off-diagonal counterpart.

\section{Derivatives and concavity of $\centroiddual$}
\label{sec:appendix-centroid}

For fixed weighted sites $(\qq, \weight)$, the energy $\centroiddual(\shift) = \energy(\dual, \combweight)$ in \eqref{eq:centroid-dual} depends on the shifts $\shift$ \emph{only} through the shifted sites $(\dual, \combweight)$ of \eqref{eq:shifted-site}.
Applying the chain rule with the Jacobian \eqref{eq:appendix-jac} and the energy gradients \eqref{eq:appendix-energy-deriv}, we obtain the first derivative of $\centroiddual$:
\begin{equation}
  \begin{aligned}
    \partial_{\shift_i} \centroiddual
      &= (\partial_{\shift_i} \dual_i) \, (\partial_{\dual_i} \energy)
         + (\partial_{\shift_i} \combweight_i) \, (\partial_{\combweight_i} \energy) \\
      &= (\partial_{\dual_i} \energy) + 2\,(\dual_i - \tcentroid_i)\, (\partial_{\combweight_i} \energy)
       = 2 \int_{\cell_i} (\tcentroid_i - \xx) \, d\xx ,
  \end{aligned}
  \label{eq:appendix-centroid-grad}
\end{equation}
thus recovering the gradient \eqref{eq:centroid-dual-gradient}.
% 
For its Hessian, we differentiate \eqref{eq:appendix-centroid-grad} once more, again via the chain rule.
Letting $\alpha, \beta$ range over the shifted-site variables $\{\dual, \combweight\}$, this yields
\begin{equation}
  \partial^2_{\shift_i \shift_j} \centroiddual
    = \sum\nolimits_{\alpha, \beta} (\partial_{\shift_i} \alpha_i)\,(\partial^2_{\alpha_i \beta_j} \energy)\,(\partial_{\shift_j} \beta_j)
      + \delta_{ij}\,(\partial^2_{\shift_i} \combweight_i)\,(\partial_{\combweight_i} \energy) ,
  \label{eq:appendix-centroid-hess-chain}
\end{equation}
with $\delta_{ij}$ indicating the Kronecker delta.
Substituting the energy Hessian \eqref{eq:appendix-energy-hess-off}--\eqref{eq:appendix-energy-hess-diag} and the Jacobian \eqref{eq:appendix-jac}, 
% and using $\smash{\xx - \tcentroid_i = (\xx - \dual_i) + (\dual_i - \tcentroid_i)}$, 
the facet terms combine into $\smash{(\xx - \tcentroid_i)(\xx - \tcentroid_j)\tr}$, recovering the facet Laplacian in \eqref{eq:centroid-hessian}.

% \paragraph{Negative semidefiniteness and null-space.}
To read off the definiteness of \eqref{eq:centroid-hessian}, we introduce the auxiliary per-cell linear field $\varphi_i(\xx) = 2\,\shift_i\tr (\xx - \tcentroid_i)$, whose variation under a shift perturbation $\delta\shift_i$ is $\delta\varphi_i(\xx) = 2\,\delta\shift_i\tr (\xx - \tcentroid_i)$.
Contracting the Hessian blocks with a perturbation $\smash{\delta\shift = \{\delta\shift_i\}}$ collapses each facet to a scalar square:
\begin{equation*}
  \delta\shift\tr \nabla^2 \centroiddual \, \delta\shift
  \!=\! 
  -\sum\nolimits_{i,j}\! \frac{1}{2\,\norm{\dual_i \!-\! \dual_j}} \int_{\face_{ij}}\! \bigl( \delta\varphi_i(\xx) - \delta\varphi_j(\xx) \bigr)^2 \, dA
  \le 0 ,
  \label{eq:appendix-centroid-nsd}
\end{equation*}
so $\nabla^2 \centroiddual \preceq 0$ and $\centroiddual$ is concave.
Being a sum of squares, this quadratic form vanishes only when $\delta\varphi_i(\xx) = \delta\varphi_j(\xx)$ on every shared facet $\face_{ij}$.
Note that this condition does \emph{not} glue the affine fields $\delta\varphi_i$ into a single global affine function: two affine functions that agree on the $(d\!-\!1)$-dimensional facet $\face_{ij}$ only agree on its supporting hyperplane, and their difference may be any multiple of the affine function defining this hyperplane.
To characterize the null-space, we name $a_i$ the affine function of the lifting in \S\ref{sec:background}, evaluated at the shifted weighted site of \eqref{eq:shifted-site}:
\begin{equation}
  a_i(\xx) = 2\,\dual_i\tr\xx - \norm{\dual_i}^2 + \combweight_i = 2\,\qq_i\tr\xx - \norm{\qq_i}^2 + 2\,\shift_i\tr(\xx - \tcentroid_i) + \weight_i ,
  \label{eq:appendix-affine}
\end{equation}
so the facet $\face_{ij}$ lies on the hyperplane $\{a_i = a_j\}$ and a perturbation induces $\delta a_i = \delta\varphi_i$ for fixed weights.
The null-space condition then reads $\smash{\delta a_i - \delta a_j = \kappa_{ij}\,(a_i - a_j)}$ for scalars $\kappa_{ij}$, \ie the perturbation keeps every facet hyperplane fixed to first order and, therefore, leaves the cells unchanged to first order.

In terms of the lifting of \S\ref{sec:background}, we have $\smash{a_i(\xx) = 2\,\dual_i\tr\xx - \sheight_i}$, where the height $\smash{\sheight_i = \norm{\dual_i}^2 - \combweight_i}$ of the shifted weighted site is related to the height $\smash{\height_i = \norm{\qq_i}^2 - \weight_i}$ of the original one by $\smash{\sheight_i = \height_i + 2\,\tcentroid_i\tr\shift_i}$.
Each shift thus moves the lifted point $\smash{(\dual_i, \sheight_i) = (\qq_i, \height_i) + (\shift_i, 2\,\tcentroid_i\tr\shift_i)}$ affinely, and \eqref{eq:appendix-affine} gives $\smash{a_i(\tcentroid_i) = 2\,\qq_i\tr\tcentroid_i - \height_i = \gamma_i}$ for every shift, with $\gamma_i$ as in \eqref{eq:pivot}.
In other words, the shift $\shift_i$ pivots the lifted hyperplane of site $i$ about the fixed point $\smash{(\tcentroid_i, \gamma_i)}$.
Applying the homotheties of \eqref{eq:lifted-homothety} to the lifted shifted sites and differentiating at $s \!=\! 1 \!+\! \mu$ with an infinitesimal offset $(\translation, -c)$ gives the infinitesimal homotheties
\begin{equation}
  \delta\dual_i = \mu\,\dual_i + \translation ,
  \qquad
  \delta\sheight_i = \mu\,\sheight_i - c ,
  \label{eq:infinitesimal-homothety}
\end{equation}
\ie $\smash{\delta a_i = \mu\,a_i + 2\,\translation\tr\xx + c}$, which always satisfy the null-space condition with $\kappa_{ij} \equiv \mu$.
As is generic \cite{Bourne2025Inverting}, we assume that they are also the only solutions, \ie that the cells determine their lifted affine functions only up to these homotheties; this fails in special configurations, such as strips of parallel cells or a boundary cell with a single neighbor, whose affine function may rotate about their shared facet.

With fixed weights, a shift perturbation changes each affine function by $\smash{\delta a_i = 2\,\delta\shift_i\tr(\xx - \tcentroid_i)}$, which vanishes at $\tcentroid_i$, so an infinitesimal homothety is realizable by the shifts exactly when it preserves all pivots.
Matching the linear and constant parts of $\delta a_i = 2\,\delta\shift_i\tr(\xx - \tcentroid_i)$ with $\mu\,a_i + 2\,\translation\tr\xx + c$ gives
\begin{equation}
  \delta\shift_i = \mu\,\dual_i + \translation ,
  \qquad
  0 = \mu\,\gamma_i + 2\,\translation\tr\tcentroid_i + c ,
  \label{eq:appendix-centroid-kernel}
\end{equation}
where the second identity evaluates the first at $\smash{\xx = \tcentroid_i}$.
Since $\delta\shift = 0$ forces $\mu = 0$ and $\translation = 0$ when the shifted sites are distinct, and then $c = 0$, the Hessian is \emph{negative definite} if and only if \eqref{eq:appendix-centroid-kernel} has only the trivial solution, \ie if and only if the vectors $\smash{(\gamma_i, \tcentroid_i, 1)}$ span $\R^{d+2}$.
Equivalently, the pivots $\smash{(\tcentroid_i, \gamma_i) \in \R^{d+1}}$ do not lie on a common hyperplane, as stated in \S\ref{sec:centroid-constraints}.
Since $n$ points in $\R^{d+1}$ always lie on a common hyperplane when $n \!\le\! d\!+\!1$, this requires $n \!\ge\! d\!+\!2$.
A vertical hyperplane ($\mu \!=\! 0$) corresponds to targets $\tcentroid_i$ lying on a common hyperplane of $\R^d$.
In the centroidal setting $\qq_i \!=\! \tcentroid_i$ with $\weight_i \!=\! 0$, the pivots $\smash{(\tcentroid_i, \norm{\tcentroid_i}^2)}$ lie on the paraboloid, and a non-vertical hyperplane ($\mu \!\neq\! 0$) cuts it along a sphere, recovering the Delaunay degeneracy of targets on a common hyperplane or sphere.
In general, $\gamma_i$ lowers the paraboloid by the power distance from $\tcentroid_i$ to its own weighted site, so the degeneracy is that of the regular triangulation of the targets weighted by these power distances.
In particular, the condition that the targets affinely span $\R^d$ is not sufficient.

\section{Derivatives and concavity of $\jointdual$}
\label{sec:appendix-joint}

The joint dual $\jointdual(\shift, \weight) = \energy(\dual, \combweight) + \sum\nolimits_i \weight_i \pvol_i$ defined in \eqref{eq:joint-dual} again depends on $(\shift, \weight)$ only through the shifted sites \eqref{eq:shifted-site}, plus the explicit term $\weight_i \pvol_i$.
Relative to \S\ref{sec:appendix-centroid}, the only new ingredient is the weight reparametrization $\partial_{\weight_i} \combweight_i = 1$ from \eqref{eq:appendix-jac}.

% \paragraph{Gradient.}
The shift gradient of $\jointdual$ is identical to \eqref{eq:appendix-centroid-grad}, while for the weights the chain rule gives
\begin{equation}
  \partial_{\weight_i} \jointdual
    = (\partial_{\weight_i} \combweight_i) \, (\partial_{\combweight_i} \energy) + \pvol_i
    = (\partial_{\combweight_i} \energy) + \pvol_i
    = -\bigl( \cellvol_i - \pvol_i \bigr) ,
  \label{eq:appendix-joint-grad}
\end{equation}
so that stacking the two gradients recovers the joint gradient \eqref{eq:joint-gradient}.

% \paragraph{Hessian.}
Differentiating a second time, the diagonal $(\shift_i,\shift_i)$ and $(\weight_i,\weight_i)$ blocks are exactly the Hessians we have already computed --- the centroid Hessian of $\centroiddual$ in \eqref{eq:centroid-hessian} and the volume Hessian of $\dualenergy$ in \eqref{eq:newton-step} --- but the joint problem also produces the mixed shift--weight blocks of \eqref{eq:joint-hessian}.
To handle all blocks at once, we introduced an auxliary variable that stacks the two multipliers of a site into $\mult_i = [\shift_i, \weight_i] \in \R^{d+1}$ and then define the affine field
\begin{equation}
  \psi_i(\xx) = 2\,\shift_i\tr (\xx - \tcentroid_i) + \weight_i = \mult_i\tr \phi_i(\xx) ,
  \qquad
  \phi_i(\xx) = \begin{bmatrix} 2(\xx - \tcentroid_i) \\ 1 \end{bmatrix} ,
  \label{eq:appendix-potential}
\end{equation}
through which $\mult_i$ enters the shifted power distance $\norm{\xx - \qq_i}^2 - \psi_i(\xx)$.
The Hessian is then the facet-based Laplacian assembled from $\phi_i$:
\begin{subequations}
  \label{eq:appendix-block}
  \begin{empheq}[left={\empheqlbrace\,}]{align}
    \partial^2_{\mult_i \mult_j} \jointdual
      &= \frac{1}{2}\frac{1}{\norm{\dual_i - \dual_j}} \int_{\face_{ij}} \phi_i(\xx)\, \phi_j(\xx)\tr \, dA,
    \label{eq:appendix-block-off}
    \\[2mm]
    \partial^2_{\mult_i \mult_i} \jointdual
      &= -\!\sum\nolimits_{j \sim i} \frac{1}{2}\frac{1}{\norm{\dual_i - \dual_j}} \int_{\face_{ij}} \phi_i(\xx)\, \phi_i(\xx)\tr \, dA.
    \label{eq:appendix-block-diag}
  \end{empheq}
\end{subequations}
Expanding $\phi_i$ splits \eqref{eq:appendix-block} into shift--shift, shift--weight, and weight--weight coefficients, with off-diagonal blocks
\begin{subequations}
  \label{eq:appendix-blocks-off}
  \begin{empheq}[left={\empheqlbrace\,}]{align}
    \partial^2_{\shift_i \shift_j\tr} \jointdual
      &= \frac{2}{\norm{\dual_i - \dual_j}} \int_{\face_{ij}}\! (\xx - \tcentroid_i)(\xx - \tcentroid_j)\tr \, dA , \label{eq:appendix-blocks-uu} \\
    \partial^2_{\shift_i \weight_j} \jointdual
      &= \frac{1}{\norm{\dual_i - \dual_j}} \int_{\face_{ij}}\! (\xx - \tcentroid_i) \, dA , \label{eq:appendix-blocks-uw} \\
    \partial^2_{\weight_i \weight_j} \jointdual
      &= \frac{1}{2}\frac{\area_{ij}}{\norm{\dual_i - \dual_j}} , \label{eq:appendix-blocks-ww}
  \end{empheq}
\end{subequations}
and diagonal blocks
\begin{subequations}
  \label{eq:appendix-blocks-diag}
  \begin{empheq}[left={\empheqlbrace\,}]{align}
    \partial^2_{\shift_i \shift_i\tr} \jointdual
      &= -\!\sum\nolimits_{j \sim i} \frac{2}{\norm{\dual_i - \dual_j}} \int_{\face_{ij}}\! (\xx - \tcentroid_i)(\xx - \tcentroid_i)\tr \, dA , \label{eq:appendix-blocks-uu-diag} \\
    \partial^2_{\shift_i \weight_i} \jointdual
      &= -\!\sum\nolimits_{j \sim i} \frac{1}{\norm{\dual_i - \dual_j}} \int_{\face_{ij}}\! (\xx - \tcentroid_i) \, dA , \label{eq:appendix-blocks-uw-diag} \\
    \partial^2_{\weight_i \weight_i} \jointdual
      &= -\!\sum\nolimits_{j \sim i} \frac{1}{2}\frac{\area_{ij}}{\norm{\dual_i - \dual_j}} , \label{eq:appendix-blocks-ww-diag}
  \end{empheq}
\end{subequations}
whose shift--shift, shift--weight, and weight--weight parts recover, respectively, the centroid Hessian \eqref{eq:centroid-hessian}, the mixed block \eqref{eq:joint-hessian}, and the volume Hessian \eqref{eq:newton-step}.

% \paragraph{Negative semidefiniteness and null space.}
Contracting the Hessian with a perturbation $\smash{\delta\mult = \{ (\delta\shift_i, \delta\weight_i) \}}$, again collapses each facet to a scalar square of the form
\begin{equation*}
  \delta\mult\tr \nabla^2 \jointdual \, \delta\mult
    = -\sum\nolimits_{i,j} \frac{1}{2\,\norm{\dual_i - \dual_j}} \int_{\face_{ij}} \bigl( \delta\psi_i(\xx) - \delta\psi_j(\xx) \bigr)^2 \, dA
    \;\le\; 0 ,
  \label{eq:appendix-joint-nsd}
\end{equation*}
now with variations of the affine field $\smash{\delta\psi_i(\xx) = 2\,\delta\shift_i\tr (\xx - \tcentroid_i) + \delta\weight_i}$.
Hence, $\nabla^2 \jointdual \preceq 0$ and $\jointdual$ is jointly concave.
As in Appendix~\ref{sec:appendix-centroid}, this quadratic form vanishes exactly when $\smash{\delta\psi_i - \delta\psi_j = \kappa_{ij}\,(a_i - a_j)}$ on every facet, now with $\smash{\delta a_i = \delta\psi_i}$ since the weights also vary.
The dilation and translation perturbation $\smash{\delta a_i = \mu\,a_i + 2\,\translation\tr\xx + c}$ always satisfies this condition and produces the only solution.
Matching the linear and constant parts as in \eqref{eq:appendix-centroid-kernel}, but now solving the constant part for the free weight perturbation, gives
\begin{equation*}
  \delta\shift_i = \mu\,\dual_i + \translation ,
  \qquad
  \delta\weight_i = \mu\,\gamma_i + 2\,\translation\tr\tcentroid_i + c ,
\end{equation*}
which is exactly the $(d\!+\!2)$-dimensional family in \eqref{eq:joint-nullspace}.
One can verify directly that it lies in the null-space, since $\smash{\delta\psi_i - \delta\psi_j = \mu\,(a_i - a_j)}$ vanishes on $\face_{ij}$.
There, $\mu$ scales the lifted shifted sites $\smash{(\dual_i, \sheight_i)}$, $\translation$ translates them horizontally, and $c$ translates them vertically, which together form the homotheties \eqref{eq:infinitesimal-homothety} that leave the power diagram unchanged (\S\ref{sec:background}).
In contrast to $\centroiddual$, these modes are always present because the free weights absorb the weight changes they require.

\section{Feasible volumes for prescribed centroids}
\label{sec:appendix-maxent}

In the centroid-only problem of \S\ref{sec:centroid-constraints}, the volumes are unknown, and the conservation identities \eqref{eq:partition-feasibility} are satisfiable only if there exist non-negative volumes $\cellvol_i$ with $\smash{\sum_i \cellvol_i \tcentroid_i = \domainVol\,\centroid_\domain}$ and $\smash{\sum_i \cellvol_i = \domainVol}$.
Drawing inspiration from \citet{Hormann2008MaxEnt}, one way to recover such volumes is to solve
\begin{subequations}
  \label{eq:maxent-feasibility}
  \begin{alignat}{2}
    & \min_{\{\cellvol_i\}} \quad && \sum\nolimits_i \cellvol_i \, \log\bigl( \cellvol_i / \domainVol \bigr) \label{eq:maxent-feasibility-obj} \\
    & \text{s.t.} \quad && \sum\nolimits_i \cellvol_i \, \tcentroid_i = \domainVol \, \centroid_\domain , \quad \sum\nolimits_i \cellvol_i = \domainVol , \label{eq:maxent-feasibility-cons}
  \end{alignat}
\end{subequations}
which is equivalent to the \emph{unconstrained} convex minimization
\begin{equation}
  \min_{\yy \in \R^{d+1}} \
    \domainVol \sum\nolimits_i \exp\bigl( -\yy\tr \lift_i \bigr)
    + \yy\tr \lift_\domain ,
  \label{eq:maxent-dual}
\end{equation}
with the homogeneous lifts $\lift_i \!=\! [\tcentroid_i, 1]$ and $\lift_\domain \!=\! \domainVol\,[\centroid_\domain, 1]$, whose minimizer recovers the volumes as $\smash{\cellvol_i = \domainVol \exp(-\yy\tr \lift_i)}$.

\section{Bordered Newton step}
\label{sec:appendix-bordered}

We detail the Newton step of Algorithm~\ref{alg:flow-constrained-newton}, which solves \eqref{eq:flow-line-constraint} by maximizing $\jointdual$ subject to the linear constraints \eqref{eq:flow-line-constraint-dil} and \eqref{eq:flow-line-constraint-trans}.

\paragraph{Constraint matrix.}
We stack the shifts and weights into $\smash{\mathbf{z} = (\shift, \weight)}$, and write the dilation constraint \eqref{eq:flow-line-constraint-dil} as the row $\smash{\mathbf{a} = (\{\pvol_i\,\dual^k_i\}_i, 0)}$ and the translation constraint \eqref{eq:flow-line-constraint-trans} along each axis $\mathbf{e}_j$ as the row $\smash{\mathbf{b}_j = (\{\pvol_i\,\mathbf{e}_j\}_i, 0)}$, both with zero entries on the weights.
We then collect these $d \!+\! 1$ rows into the auxiliary matrix
\begin{equation}
  \Cmat = \begin{bmatrix} \mathbf{a} & \mathbf{b}_1 & \cdots & \mathbf{b}_d \end{bmatrix}\tr ,
  \label{eq:appendix-constraint-matrix}
\end{equation}
so that the constraints \eqref{eq:flow-line-constraint-dil} and \eqref{eq:flow-line-constraint-trans} read $\Cmat\,\mathbf{z} = 0$.

\paragraph{Newton system.}
Denoting by $\smash{\mathbf{g} = \nabla\jointdual(\mathbf{z})}$ the gradient \eqref{eq:joint-gradient}, the constraint multipliers $\smash{\boldsymbol{\nu} = 2\,(s, t_1, \dots, t_d) \in \R^{d+1}}$, 
% with the translation $\smash{\translation = (t_1, \dots, t_d)}$, and the Newton step $\delta\mathbf{z}$, 
we solve the bordered system
\begin{equation}
  \begin{bmatrix} \Nmat & \Cmat\tr \\ \Cmat & 0 \end{bmatrix}
  \begin{bmatrix} \delta\mathbf{z} \\ \boldsymbol{\nu} \end{bmatrix}
  =
  \begin{bmatrix} \mathbf{g} \\ 0 \end{bmatrix},
  \label{eq:flow-constrained-newton}
\end{equation}
%
where the matrix
\begin{equation}
  \Nmat = -\nabla^2\jointdual(\mathbf{z}) + \varepsilon\,\mathrm{diag}(\pvol)
  \label{eq:appendix-bordered-matrix}
\end{equation}
is the negated Hessian of $\jointdual$ perturbed by a small $\varepsilon \!>\! 0$ multiple of the target volumes $\pvol$ to remove its null-space \eqref{eq:joint-nullspace}.
Since the negated Hessian is positive semi-definite, $\Nmat$ is positive definite and sparse, with the same pattern as the Hessian, so it admits a sparse Cholesky factorization.

\paragraph{Block elimination.}
Eliminating $\delta\mathbf{z}$ from the first row of \eqref{eq:flow-constrained-newton} gives $\smash{\delta\mathbf{z} = \Nmat^{-1} \bigl( \mathbf{g} - \Cmat\tr \boldsymbol{\nu} \bigr)}$, which we compute from the $d \!+\! 2$ sparse solves
\begin{subequations}
  \label{eq:flow-constrained-elimination}
  \begin{equation}
    \Nmat\,\mathbf{z}_0 = \mathbf{g} ,
    \qquad
    \Nmat\,\mathbf{z}_a = 2\,\mathbf{a} ,
    \qquad
    \Nmat\,\mathbf{z}_{b_j} = 2\,\mathbf{b}_j \ \ (j = 1, \dots, d) ,
    \label{eq:flow-constrained-elimination-solve}
  \end{equation}
  which share a single factorization of $\Nmat$.
  Substituting into the last row $\Cmat\,\delta\mathbf{z} = 0$ yields the dense $(d \!+\! 1) \times (d \!+\! 1)$ normal equation:
  \begin{equation}
    \begin{bmatrix}
      \mathbf{a}\tr\mathbf{z}_a & \mathbf{a}\tr\mathbf{z}_{b_1} & \cdots & \mathbf{a}\tr\mathbf{z}_{b_d} \\
      \mathbf{b}_1\tr\mathbf{z}_a & \mathbf{b}_1\tr\mathbf{z}_{b_1} & \cdots & \mathbf{b}_1\tr\mathbf{z}_{b_d} \\
      \vdots & \vdots & \ddots & \vdots \\
      \mathbf{b}_d\tr\mathbf{z}_a & \mathbf{b}_d\tr\mathbf{z}_{b_1} & \cdots & \mathbf{b}_d\tr\mathbf{z}_{b_d}
    \end{bmatrix}
    \begin{bmatrix} s \\ t_1 \\ \vdots \\ t_d \end{bmatrix}
    =
    \begin{bmatrix} \mathbf{a}\tr\mathbf{z}_0 \\ \mathbf{b}_1\tr\mathbf{z}_0 \\ \vdots \\ \mathbf{b}_d\tr\mathbf{z}_0 \end{bmatrix} ,
    \label{eq:flow-constrained-elimination-gauge}
  \end{equation}
  which is symmetric positive definite.
  The Newton step follows as
  \begin{equation}
    \delta\mathbf{z} = \mathbf{z}_0 - s\,\mathbf{z}_a - \sum\nolimits_j t_j\,\mathbf{z}_{b_j}.
    \label{eq:flow-constrained-elimination-update}
  \end{equation}
\end{subequations}

\section{Algorithms}
\label{sec:appendix-algorithms}

We summarize the time steps of the constrained (\S\ref{sec:flow-constrained}) and penalty-based (\S\ref{sec:flow-penalty}) integrators, together with the Newton solves they rely on.
Both solves perturb the negated Hessian by $\smash{\varepsilon\,\mathrm{diag}(\pvol)}$, and backtrack the step size $\omega$ until the objective increases.

\begin{algorithm}[h]
  \caption{Time step of the constrained integrator (\S\ref{sec:flow-constrained}).}
  \label{alg:flow-constrained}
  \begin{algorithmic}[1]
    \Require particles $\qq^k$, velocities $\vel^k$, generators $\dual^k$ and shifted weights $\combweight^k$ of the power diagram with centroids $\qq^k$, volumes $\pvol$, time step $\tstep$, strength $\stiff$
    \State $\pred^{k+1} \gets (1 + \tstep^2/\stiff^2)\,\qq^k + \tstep\,\vel^k$ \Comment{prediction \eqref{eq:flow-shake-step-pred}}
    \State $(\shift, \weight, s, \translation) \gets$ Alg.~\ref{alg:flow-constrained-newton} with $\qq^k$, $\pred^{k+1}$, $\dual^k$, $\combweight^k$
    \State $\qq^{k+1} \gets \pred^{k+1} - s\,\dual^k - \translation$ \Comment{correction \eqref{eq:flow-shake-step-corr}}
    \State $\vel^{k+1} \gets (\qq^{k+1} - \qq^k) / \tstep$
    \State $\dual^{k+1} \gets \qq^k + \shift$ \Comment{\eqref{eq:shifted-site-pos}}
    \State $\combweight^{k+1} \gets \weight + \norm{\shift}^2 + 2\,\shift\tr (\qq^k - \pred^{k+1})$ \Comment{\eqref{eq:shifted-site-weight}}
  \end{algorithmic}
\end{algorithm}

\begin{algorithm}[h]
  \caption{Bordered Newton solve of \eqref{eq:flow-line-constraint} (Appendix~\ref{sec:appendix-bordered}).}
  \label{alg:flow-constrained-newton}
  \begin{algorithmic}[1]
    \Require sites $\qq$, targets $\pred$, direction $\dual^k$, initial weights $\weight^0$, volumes $\pvol$, regularization $\varepsilon \!>\! 0$, tolerance $\tau$
    \State build the constraint matrix $\Cmat = [\mathbf{a}, \mathbf{b}_1, \dots, \mathbf{b}_d]\tr$ \Comment{\eqref{eq:appendix-constraint-matrix}}
    \State $\mathbf{z} = (\shift, \weight) \gets (0, \weight^0)$
    \Repeat
      \State build the power diagram of the shifted sites \eqref{eq:shifted-site}
      \State $\mathbf{g} \gets \nabla\jointdual(\mathbf{z})$; \ $\Nmat \gets$ \eqref{eq:appendix-bordered-matrix}
      \State solve for $\mathbf{z}_0$, $\mathbf{z}_a$, $\mathbf{z}_{b_j}$ \Comment{$d \!+\! 2$ sparse solves \eqref{eq:flow-constrained-elimination-solve}}
      \State solve for the gauge $(s, \translation)$ \Comment{dense system \eqref{eq:flow-constrained-elimination-gauge}}
      \State $\delta\mathbf{z} \gets \mathbf{z}_0 - s\,\mathbf{z}_a - \sum_j t_j\,\mathbf{z}_{b_j}$ \Comment{\eqref{eq:flow-constrained-elimination-update}}
      \State $\mathbf{z} \gets \mathbf{z} + \omega\,\delta\mathbf{z}$, halving $\omega$ from $1$ until $\jointdual$ increases
    \Until{$\norm{\mathbf{g} - \Cmat\tr \boldsymbol{\nu}} < \tau$, with $\boldsymbol{\nu} = 2\,(s, t_1, \dots, t_d)$} \Comment{stationarity of \eqref{eq:flow-constrained-newton}}
    \State \Return $(\shift, \weight)$ and the gauge $(s, \translation)$
  \end{algorithmic}
\end{algorithm}

\begin{algorithm}[h]
  \caption{Time step of the penalty-based integrator (\S\ref{sec:flow-penalty}).}
  \label{alg:flow-penalty}
  \begin{algorithmic}[1]
    \Require particles $\qq^k$, velocities $\vel^k$, maximizer $(\shift^k, \weight^k)$ of $\jointdual_\alpha$ with sites and targets $\qq^k$, volumes $\pvol$, time step $\tstep$, strength $\stiff$, stiffness $\alpha$
    \State $\vel^{k+1} \gets \vel^k - (1 + 1/\alpha)\,(\tstep/\stiff^2)\,\shift^k$ \Comment{velocity update \eqref{eq:flow-softened-step-vel}}
    \State $\qq^{k+1} \gets \qq^k + \tstep\,\vel^{k+1}$ \Comment{position update \eqref{eq:flow-softened-step-pos}}
    \State $(\shift^{k+1}, \weight^{k+1}) \gets$ Alg.~\ref{alg:newton} with sites and targets $\qq^{k+1}$, weights $\weight^k$
  \end{algorithmic}
\end{algorithm}

\begin{algorithm}[h]
  \caption{Newton solve of $\jointdual_\alpha$ (or of $\jointdual$, for $\alpha = \infty$).}
  \label{alg:newton}
  \begin{algorithmic}[1]
    \Require sites $\qq$, targets $\tcentroid$, initial weights $\weight^0$, volumes $\pvol$, stiffness $\alpha$ (optional), regularization $\varepsilon \!>\! 0$, tolerance $\tau$
    \State $\mathbf{z} = (\shift, \weight) \gets (0, \weight^0)$
    \Repeat
      \State build the power diagram of the shifted sites \eqref{eq:shifted-site}
      \State $\mathbf{g} \gets \nabla\jointdual_\alpha(\mathbf{z})$; \ $\Nmat \gets -\nabla^2\jointdual_\alpha(\mathbf{z}) + \varepsilon\,\mathrm{diag}(\pvol)$ \Comment{as in \eqref{eq:appendix-bordered-matrix}}
      \State solve $\Nmat\,\delta\mathbf{z} = \mathbf{g}$
      \State $\mathbf{z} \gets \mathbf{z} + \omega\,\delta\mathbf{z}$, halving $\omega$ from $1$ until $\jointdual_\alpha$ increases
    \Until{$\norm{\mathbf{g}} < \tau$}
    \State \Return $(\shift, \weight)$
  \end{algorithmic}
\end{algorithm}

\end{document}
